\chapter{LCS con un numero di simboli rossi multiplo di 3}\section{Definizione del Problema}Siano date due sequenze:
\begin{align*}
X &= \langle x_1, x_2, \dots, x_m \rangle \
Y &= \langle y_1, y_2, \dots, y_n \rangle
\end{align*}
su un generico alfabeto finito $\Sigma$, e sia definita una funzione colore:
$$\text{col} : \Sigma \to \{\text{rosso}, \text{non-rosso}\}$$
Il problema richiede di trovare la più lunga sottosequenza comune (LCS) tra $X$ e $Y$ tale che il numero complessivo di elementi di colore rosso sia un \textbf{multiplo di 3} (ovvero pari a $0, 3, 6, 9, \dots$).\section{Definizione del Generico Sottoproblema e del Coefficiente}Per modellare il vincolo di divisibilità per 3 mediante la programmazione dinamica, introduciamo uno stato modulare:
$$p \in \{0, 1, 2\}$$
che rappresenta il resto della divisione intera per 3 del numero di simboli rossi inclusi nella sottosequenza comune (aritmetica modulare $\pmod 3$):
\begin{itemize}
\item $p = 0$: il numero di rossi è $3k$ ($0 \pmod 3$, cioè multiplo di 3);
\item $p = 1$: il numero di rossi è $3k + 1$ ($1 \pmod 3$);
\item $p = 2$: il numero di rossi è $3k + 2$ ($2 \pmod 3$).
\end{itemize}\subsection{Sottoproblema di dimensione $(i, j, p)$}
Dati i prefissi $X_i = \langle x_1, \dots, x_i \rangle$ e $Y_j = \langle y_1, \dots, y_j \rangle$ con:
$$i \in \{0, 1, \dots, m\}, \quad j \in \{0, 1, \dots, n\}, \quad p \in \{0, 1, 2\}$$
il sottoproblema consiste nel calcolare una più lunga sottosequenza comune tra $X_i$ e $Y_j$ che contenga un numero di simboli rossi congruo a $p \pmod 3$.\subsection{Definizione del Coefficiente $c_{i,j}^p$}
Definiamo il coefficiente $c_{i,j}^p$ come:
$$c_{i,j}^p = \begin{cases} \vert{}Z^{(i,j,p)}\vert{} & \text{se esiste una sottosequenza comune tra } X_i \text{ e } Y_j \text{ con numero di rossi } \equiv p \pmod 3 \\ -\infty & \text{altrimenti (configurazione inammissibile)} \end{cases}$$\subsection{Soluzione Ottima del Problema Principale}
Il problema principale corrisponde all'intera estensione delle due sequenze con richiesta di un numero di rossi multiplo di 3 ($p = 0$):
\begin{itemize}
\item \textbf{Lunghezza ottima:} $c_{m,n}^0$ (se $c_{m,n}^0 = -\infty$, non esiste alcuna sottosequenza ammissibile);
\item \textbf{Sottosequenza ottima:} $Z^{(m,n,0)}$.
\end{itemize}\section{Equazioni di Ricorrenza}\subsection{Casi Base}
I casi base si hanno quando almeno uno dei due prefissi è vuoto, ossia quando $i = 0 \lor j = 0$:
$$\forall p \in \{0, 1, 2\}, \quad (i = 0 \lor j = 0) \implies c_{i,j}^p = \begin{cases} 0 & \text{se } p = 0 \quad (\text{la sequenza vuota } \langle \rangle \text{ ha } 0 \text{ rossi, e } 0 \equiv 0 \pmod 3) \\ -\infty & \text{se } p \in \{1, 2\} \quad (\text{impossibile avere resti positivi con } \langle \rangle) \end{cases}$$\subsection{Passo Ricorsivo ($i > 0 \land j > 0$)}
Siano $i \in \{1, \dots, m\}$ e $j \in \{1, \dots, n\}$. Si distinguono i seguenti casi:\paragraph{Caso 1: $x_i \neq y_j$}
Gli elementi finali non coincidono; il miglior allineamento non può includerli contemporaneamente:
$$\forall p \in \{0, 1, 2\}, \quad c_{i,j}^p = \max \{ c_{i-1,j}^p, \, c_{i,j-1}^p \}$$\paragraph{Caso 2: $x_i = y_j$ e $\text{col}(x_i) \neq \text{rosso}$}
Il carattere comune non è rosso; la sua inclusione incrementa la lunghezza di 1 senza variare il numero di rossi accumulati:
$$\forall p \in \{0, 1, 2\}, \quad c_{i,j}^p = c_{i-1,j-1}^p + 1$$
\textit{Nota convenzionale:} se $c_{i-1,j-1}^p = -\infty$, si assume $-\infty + 1 = -\infty$.\paragraph{Caso 3: $x_i = y_j$ e $\text{col}(x_i) = \text{rosso}$}
Il carattere comune è rosso. Includendolo, la parità modulare passa da $(p - 1) \pmod 3$ a $p$:
\begin{itemize}
\item \textbf{Se $p = 0$:} possiamo scegliere se includere $x_i$ (agganciandoci a una soluzione ottima con resto $2$, poiché $(2 + 1) \equiv 0 \pmod 3$) oppure scartarlo mantenendo il resto $0$:
$$c_{i,j}^0 = \max \{ c_{i-1,j-1}^2 + 1, \, c_{i-1,j-1}^0 \}$$
\item \textbf{Se $p = 1$:} per ottenere resto $1$ dobbiamo necessariamente accodare $x_i$ a una soluzione con resto $0$ (infatti $0 + 1 = 1$):
$$c_{i,j}^1 = c_{i-1,j-1}^0 + 1$$
\item \textbf{Se $p = 2$:} per ottenere resto $2$ dobbiamo necessariamente accodare $x_i$ a una soluzione con resto $1$ (infatti $1 + 1 = 2$):
$$c_{i,j}^2 = c_{i-1,j-1}^1 + 1$$
\end{itemize}\section{Algoritmo Iterativo Bottom-Up}Riportiamo lo pseudocodice corretto per la costruzione della tabella tridimensionale $c[0\dots m, 0\dots n, 0\dots 2]$.\begin{algorithm}[H]
\caption{Calcolo Bottom-Up per LCS con rossi multipli di 3}
\begin{algorithmic}[1]
\Procedure{LCS-M3R}{$X, Y$}
\State $m \gets \vert{}X\vert{}$
\State $n \gets \vert{}Y\vert{}$
\State Inizializza matrice $c[0 \dots m, \, 0 \dots n, \, 0 \dots 2]$ con $-\infty$\Comment{\textbf{Inizializzazione dei Casi Base}}
\For{$i \gets 0 \textbf{ to } m$}
    \State $c[i, 0, 0] \gets 0$
    \State $c[i, 0, 1] \gets -\infty$
    \State $c[i, 0, 2] \gets -\infty$
\EndFor
\For{$j \gets 0 \textbf{ to } n$}
    \State $c[0, j, 0] \gets 0$
    \State $c[0, j, 1] \gets -\infty$
    \State $c[0, j, 2] \gets -\infty$
\EndFor

\Comment{\textbf{Riempimento Bottom-Up}}
\For{$i \gets 1 \textbf{ to } m$}
    \For{$j \gets 1 \textbf{ to } n$}
        \If{$x_i \neq y_j$}
            \For{$p \gets 0 \textbf{ to } 2$}
                \State $c[i, j, p] \gets \max\{ c[i-1, j, p], \, c[i, j-1, p] \}$
            \EndFor
        \Else \Comment{$x_i = y_j$}
            \If{$\text{col}(x_i) \neq \text{rosso}$}
                \For{$p \gets 0 \textbf{ to } 2$}
                    \State $c[i, j, p] \gets c[i-1, j-1, p] + 1$
                \EndFor
            \Else \Comment{$x_i = y_j \land \text{col}(x_i) = \text{rosso}$}
                \State $c[i, j, 0] \gets \max\{ c[i-1, j-1, 2] + 1, \, c[i-1, j-1, 0] \}$
                \State $c[i, j, 1] \gets c[i-1, j-1, 0] + 1$
                \State $c[i, j, 2] \gets c[i-1, j-1, 1] + 1$
            \EndIf
        \EndIf
    \EndFor
\EndFor
\State \Return $c$
\EndProcedure
\end{algorithmic}
\end{algorithm}\section{Algoritmo di Ricostruzione della Soluzione}La procedura di ricostruzione risale la tabella $c$ determinando all'indietro le scelte ottime effettuate. La stampa dei caratteri avviene in post-ordine (dopo la chiamata ricorsiva), garantendo che la sottosequenza sia stampata da sinistra verso destra.\begin{algorithm}[H]
\caption{Ricostruzione Ricorsiva della LCS vincolata}
\begin{algorithmic}[1]
\Procedure{PRINT-LCS-M3R}{$c, X, Y, i, j, p$}
\If{$i = 0 \lor j = 0$}
\State \Return
\EndIf\If{$x_i \neq y_j$}
    \If{$c[i, j, p] == c[i-1, j, p]$}
        \State \Call{PRINT-LCS-M3R}{$c, X, Y, i-1, j, p$}
    \Else
        \State \Call{PRINT-LCS-M3R}{$c, X, Y, i, j-1, p$}
    \EndIf
\Else \Comment{$x_i = y_j$}
    \If{$\text{col}(x_i) \neq \text{rosso}$}
        \State \Call{PRINT-LCS-M3R}{$c, X, Y, i-1, j-1, p$}
        \State \textbf{print} $x_i$
    \Else \Comment{$\text{col}(x_i) = \text{rosso}$}
        \If{$p = 0$}
            \If{$c[i, j, 0] == c[i-1, j-1, 2] + 1$}
                \State \Call{PRINT-LCS-M3R}{$c, X, Y, i-1, j-1, 2$}
                \State \textbf{print} $x_i$
            \Else
                \State \Call{PRINT-LCS-M3R}{$c, X, Y, i-1, j-1, 0$}
            \EndIf
        \ElsIf{$p = 1$}
            \State \Call{PRINT-LCS-M3R}{$c, X, Y, i-1, j-1, 0$}
            \State \textbf{print} $x_i$
        \ElsIf{$p = 2$}
            \State \Call{PRINT-LCS-M3R}{$c, X, Y, i-1, j-1, 1$}
            \State \textbf{print} $x_i$
        \EndIf
    \EndIf
\EndIf
\EndProcedure
\end{algorithmic}
\end{algorithm}\subsection{Procedura Principale (Driver)}
Per avviare la ricostruzione sulla soluzione globale del problema:
\begin{algorithmic}[1]
\If{$c[m, n, 0] == -\infty$}
\State \textbf{print} "Nessuna sottosequenza comune con numero di rossi multiplo di 3"
\Else
\State \Call{PRINT-LCS-M3R}{$c, X, Y, m, n, 0$}
\EndIf
\end{algorithmic}\section{Analisi della Complessità}\subsection{Complessità Temporale}
\begin{itemize}
\item \textbf{Calcolo Bottom-Up:} I cicli su $i$ e $j$ iterano $(m+1)(n+1)$ volte. All'interno, lo stato $p$ assume esattamente 3 valori ($0, 1, 2$). Ogni operazione interna (confronto o incremento) richiede tempo $O(1)$. Il tempo totale è quindi:
$$T(m, n) = \Theta(m \cdot n \cdot 3) = \Theta(m \cdot n)$$
\item \textbf{Ricostruzione:} Ad ogni passo ricorsivo, almeno uno tra gli indici $i$ e $j$ decresce di almeno 1. La procedura esegue al più $m + n$ passi, ciascuno di costo $O(1)$:
$$T_{\text{ricostruzione}}(m, n) = O(m + n)$$
\end{itemize}\subsection{Complessità Spaziale}
L'algoritmo memorizza la tabella $c$ di dimensioni $(m + 1) \times (n + 1) \times 3$:
$$S(m, n) = \Theta(m \cdot n)$$
Lo spazio ausiliario per lo stack della ricorsione nella procedura di ricostruzione è limitato a $O(m + n)$.